% \documentclass[11pt,a4paper]{article}

% \usepackage[margin=1in]{geometry}
% \usepackage[T1]{fontenc}
% \usepackage{lmodern}
% \usepackage{amsmath,amssymb,amsfonts}
% \usepackage{graphicx}
% \usepackage{booktabs}
% \usepackage{multirow}
% \usepackage{hyperref}
% \usepackage{xcolor}
% \usepackage{float}
% \usepackage{caption}
% \usepackage{subcaption}
% \usepackage{cite}
% \usepackage{titlesec}

% \hypersetup{
%     colorlinks=true,
%     linkcolor=blue!70!black,
%     citecolor=green!50!black,
%     urlcolor=blue!80!black
% }

% \titleformat{\section}{\large\bfseries}{\thesection}{1em}{}
% \titleformat{\subsection}{\normalsize\bfseries}{\thesubsection}{1em}{}

% \title{
%     \vspace{-1.5cm}
%     \textbf{\Large EN3150: Pattern Recognition \& Machine Intelligence}\\
%     \textbf{\large Assignment 03: Resource-Constrained CNN for Edge Image Classification}\\
%     \vspace{0.3cm}
%     \large Empirical Evaluation of Depthwise Separable Convolutions, Optimizer Dynamics, and Lightweight SOTA Transfer Learning
% }

% \author{
%     \textbf{Group Members:} \\
%     Thevinu Fernando \quad Januka Sanujaya \quad Thisura Samuditha \quad Chamod Dharmakeerthi \\
%     \textit{Department of Electronic \& Telecommunication Engineering} \\
%     \textit{University of Moratuwa} \\
%     \texttt{GitHub: https://github.com/thevinufernando/cnn-based-vegetable-classification-model}
% }

% \date{\today}

% \begin{document}

% \maketitle

% \begin{abstract}
% Deploying deep convolutional neural networks (CNNs) on resource-constrained edge hardware (e.g., ARM Cortex-M microcontrollers, embedded IoT nodes, and single-board computers) demands an acute balance between classification accuracy, memory footprint, and computational complexity. In this assignment, we investigate edge-oriented image classification across 15 vegetable categories downscaled to $64 \times 64$ pixels. We design, implement, and benchmark two custom architectures from scratch: \textbf{Model A} (Standard CNN, 259,631 trainable parameters) and \textbf{Model B} (Depthwise Separable CNN, 40,175 trainable parameters, adhering strictly to a sub-100k budget). Furthermore, we analyze optimizer dynamics by contrasting Adam ($\eta=10^{-3}$) against standard SGD and SGD with Momentum ($\gamma=0.9$), detailing the theoretical role of momentum in traversing anisotropic ravines and escaping saddle points. Finally, we fine-tune two mobile state-of-the-art architectures (\textbf{MobileNetV2} and \textbf{EfficientNet-B0}) under identical data splits. Our empirical findings demonstrate that Model B achieves \textbf{97.70\% test accuracy}—outperforming Model A (\textbf{97.57\%}) while reducing parameters by \textbf{84.5\%}, disk footprint by \textbf{80.0\%}, and training time per epoch by \textbf{47.4\%}. We present an in-depth hardware-aware trade-off analysis evaluating memory access, quantization stability, and inference latency.
% \end{abstract}

% \section{Resource-Constrained CNN for Edge Image Classification}

% \subsection{Dataset Selection \& Low-Resolution Edge Sensor Emulation}
% Edge devices, such as agricultural field sensors, robotic harvesters, and low-power smart cameras, operate under stringent energy and static RAM (SRAM) constraints. For this study, we selected the \textbf{Vegetable Image Dataset} (from Kaggle), encompassing 21,000 RGB images across 15 distinct classes: \textit{Bean, Bitter Gourd, Bottle Gourd, Brinjal, Broccoli, Cabbage, Capsicum, Carrot, Cauliflower, Cucumber, Papaya, Potato, Pumpkin, Radish, and Tomato}.

% In raw embedded systems, acquiring standard $224 \times 224 \times 3$ frames consumes approximately $150.5\text{ KB}$ per uncompressed image buffer. On microcontrollers featuring only $256\text{ KB}$ to $512\text{ KB}$ of total SRAM, buffering high-resolution images leaves insufficient memory for intermediate feature maps and activation buffers. Consequently, downscaling images to a maximum resolution of \textbf{$64 \times 64$ pixels} compresses the raw uncompressed buffer down to \textbf{$12.288\text{ KB}$} ($64 \times 64 \times 3$ bytes). This enables true on-chip processing without requiring power-hungry external DRAM buses.

% \subsection{Data Partitioning \& Leakage Prevention}
% To prevent data contamination between training and evaluation phases, MD5 checksum deduplication was performed, removing identical duplicate images. The dataset was partitioned into a strict stratified split:
% \begin{itemize}
%     \item \textbf{Training Set (70\%)}: 14,996 images
%     \item \textbf{Validation Set (15\%)}: 3,000 images
%     \item \textbf{Testing Set (15\%)}: 3,000 images
%     \item \textbf{Total Unique Images}: 20,996 images
% \end{itemize}
% This fixed split is deterministically recorded in \texttt{splits.csv} to guarantee absolute reproducibility across all experiments.

% \section{Custom Architecture Design [20 Marks]}

% \subsection{Theoretical Foundations: Standard vs. Depthwise Separable Convolutions}
% A standard 2D convolution transforms an input feature map $F \in \mathbb{R}^{D_F \times D_F \times M}$ into an output feature map $G \in \mathbb{R}^{D_F \times D_F \times N}$ using $N$ filters of dimension $D_K \times D_K \times M$. Assuming unit stride and same-padding, the total Multiply-Accumulate (MAC) computational cost is:
% \begin{equation}
%     \text{Cost}_{\text{Standard}} = D_K \times D_K \times M \times N \times D_F \times D_F
% \end{equation}

% In contrast, \textbf{Depthwise Separable Convolutions} (pioneered in MobileNets and Xception) factorize this joint spatial-channel operation into two decoupled stages:
% \begin{enumerate}
%     \item \textbf{Depthwise Convolution}: Applies a single $D_K \times D_K \times 1$ spatial filter per input channel:
%     \begin{equation}
%         \text{Cost}_{\text{Depthwise}} = D_K \times D_K \times M \times D_F \times D_F
%     \end{equation}
%     \item \textbf{Pointwise Convolution}: A $1 \times 1$ convolution combining linear projections across all $M$ channels to produce $N$ channels:
%     \begin{equation}
%         \text{Cost}_{\text{Pointwise}} = 1 \times 1 \times M \times N \times D_F \times D_F
%     \end{equation}
% \end{enumerate}

% Summing both stages, the total computational complexity becomes:
% \begin{equation}
%     \text{Cost}_{\text{Separable}} = D_K^2 \cdot M \cdot D_F^2 + M \cdot N \cdot D_F^2
% \end{equation}
% The relative computational and parameter reduction ratio is given by:
% \begin{equation}
%     \frac{\text{Cost}_{\text{Separable}}}{\text{Cost}_{\text{Standard}}} = \frac{D_K^2 \cdot M \cdot D_F^2 + M \cdot N \cdot D_F^2}{D_K^2 \cdot M \cdot N \cdot D_F^2} = \frac{1}{N} + \frac{1}{D_K^2}
% \end{equation}
% For standard $3 \times 3$ kernels ($D_K = 3$) and typical channel depths ($N \ge 32$):
% \begin{equation}
%     \frac{1}{N} + \frac{1}{9} \approx \frac{1}{9} \approx 11.1\%
% \end{equation}
% Depthwise separable convolutions achieve an \textbf{$\sim 8\times$ to $9\times$ reduction in MAC operations} ($\sim 88.9\%$ computational savings) compared to dense standard convolutions.

% \subsection{Architectural Specifications \& Parameter Verification}

% \subsubsection{Model A: Standard CNN}
% Model A interleaves four standard 2D convolution blocks (with Batch Normalization and ReLU), $2 \times 2$ Max Pooling, Global Average Pooling, and a Dense classification head:
% \begin{itemize}
%     \item \textbf{Block 1}: Conv2D(32, $3 \times 3$) $\rightarrow$ BN $\rightarrow$ ReLU $\rightarrow$ MaxPool($2 \times 2$) $\rightarrow$ Output: $32 \times 32 \times 32$. Params: 864.
%     \item \textbf{Block 2}: Conv2D(64, $3 \times 3$) $\rightarrow$ BN $\rightarrow$ ReLU $\rightarrow$ MaxPool($2 \times 2$) $\rightarrow$ Output: $16 \times 16 \times 64$. Params: 18,432.
%     \item \textbf{Block 3}: Conv2D(128, $3 \times 3$) $\rightarrow$ BN $\rightarrow$ ReLU $\rightarrow$ MaxPool($2 \times 2$) $\rightarrow$ Output: $8 \times 8 \times 128$. Params: 73,728.
%     \item \textbf{Block 4}: Conv2D(128, $3 \times 3$) $\rightarrow$ BN $\rightarrow$ ReLU $\rightarrow$ MaxPool($2 \times 2$) $\rightarrow$ Output: $4 \times 4 \times 128$. Params: 147,456.
%     \item \textbf{Head}: GlobalAveragePooling2D $\rightarrow$ Dense(128, ReLU) $\rightarrow$ Dropout(0.3) $\rightarrow$ Dense(15, Softmax). Params: 18,447.
%     \item \textbf{Total Parameters}: \textbf{260,335} (Trainable: \textbf{259,631}).
% \end{itemize}

% \subsubsection{Model B: Lightweight Separable CNN (Sub-100k Budget)}
% Model B enforces strict parameter efficiency by utilizing depthwise separable blocks:
% \begin{itemize}
%     \item \textbf{Entry Conv}: Conv2D(16, $3 \times 3$) $\rightarrow$ BN $\rightarrow$ ReLU6 $\rightarrow$ MaxPool($2 \times 2$) $\rightarrow$ Output: $32 \times 32 \times 16$. Params: 496.
%     \item \textbf{Separable Block 1}: DepthwiseConv2D($3 \times 3$) $\rightarrow$ PointwiseConv2D(32, $1 \times 1$) $\rightarrow$ MaxPool($2 \times 2$). Params: 848.
%     \item \textbf{Separable Block 2}: DepthwiseConv2D($3 \times 3$) $\rightarrow$ PointwiseConv2D(64, $1 \times 1$) $\rightarrow$ MaxPool($2 \times 2$). Params: 2,720.
%     \item \textbf{Separable Block 3}: DepthwiseConv2D($3 \times 3$) $\rightarrow$ PointwiseConv2D(128, $1 \times 1$) $\rightarrow$ MaxPool($2 \times 2$). Params: 9,536.
%     \item \textbf{Separable Block 4}: DepthwiseConv2D($3 \times 3$) $\rightarrow$ PointwiseConv2D(128, $1 \times 1$). Params: 18,560.
%     \item \textbf{Head}: GlobalAveragePooling2D $\rightarrow$ Dense(64, ReLU) $\rightarrow$ Dropout(0.2) $\rightarrow$ Dense(15, Softmax). Params: 9,231.
%     \item \textbf{Total Parameters}: \textbf{41,391} (Trainable: \textbf{40,175}).
% \end{itemize}

% \noindent\textbf{Constraint Compliance}: Model B strictly satisfies the assignment criterion ($\le 100,000$ trainable parameters) by containing exactly \textbf{40,175 parameters} ($\mathbf{59.8\%}$ below the allowable ceiling).

% \subsection{Hardware-Aware Activation Function Justification}
% In resource-constrained microcontrollers and edge digital signal processors (DSPs):
% \begin{enumerate}
%     \item \textbf{Elimination of Transcendental Exponentials}: Activations such as Sigmoid ($\sigma(x) = \frac{1}{1 + e^{-x}}$) and Tanh require floating-point exponential evaluations ($e^x$). Edge microprocessors (e.g., ARM Cortex-M0+/M3/M4) lack hardware transcendental math units. Evaluating $e^x$ necessitates Taylor series expansions or CORDIC algorithms that consume 50–150 clock cycles per element, or lookup tables (LUTs) that deplete scarce SRAM.
%     \item \textbf{Single-Cycle Execution of ReLU}: Standard ReLU ($f(x) = \max(0, x)$) evaluates in a single CPU cycle via a conditional branch or bitwise clamp instruction (\texttt{VMAX} in ARM NEON/CMSIS-NN).
%     \item \textbf{Quantization Stability via ReLU6}: In Model B, we employ \textbf{ReLU6} ($f(x) = \min(\max(0, x), 6)$). Edge accelerators deploy fixed-point 8-bit integer quantization (INT8) to achieve $4\times$ memory compression. Unbounded activations expand the dynamic range $[0, x_{\max}]$, reducing the numerical resolution of small activations when mapped to 256 discrete integer bins. Capping activation at 6.0 guarantees a fixed, high-resolution INT8 step size ($\Delta = 6.0 / 256 \approx 0.0234$), preventing numerical overflow and saturation.
% \end{enumerate}

% \section{Optimizer Selection \& Tuning [15 Marks]}

% \subsection{Selected Optimizer: Adam}
% We selected \textbf{Adam (Adaptive Moment Estimation)} with an initial learning rate $\eta = 0.001$, $\beta_1 = 0.9$, $\beta_2 = 0.999$, and $\epsilon = 10^{-7}$. 

% \noindent\textbf{Justification}: In deep vision architectures, gradients with respect to early convolutional filters are substantially smaller than gradients at the classification head. Adam computes individual adaptive learning rates for each parameter by maintaining exponentially decaying moving averages of past gradients ($m_t$, first moment) and uncentered variance ($v_t$, second moment). This avoids vanishing updates in early layers while preventing destructive updates in dense layers.

% \subsection{Comparative Analysis: Standard SGD vs. SGD with Momentum}
% Model B was trained under identical data splits across three optimizer configurations:
% \begin{enumerate}
%     \item \textbf{Adam}: $\eta = 0.001$
%     \item \textbf{Standard SGD}: $\theta_{t+1} = \theta_t - \eta \nabla L(\theta_t)$, $\eta = 0.01$, momentum $= 0.0$
%     \item \textbf{SGD with Momentum}: $v_{t+1} = \gamma v_t + \eta \nabla L(\theta_t)$, $\theta_{t+1} = \theta_t - v_{t+1}$, with $\gamma = 0.9$.
% \end{enumerate}

% \begin{figure}[H]
%     \centering
%     \includegraphics[width=0.92\linewidth]{../results/figures/optimizer_comparison.png}
%     \caption{Loss convergence and validation accuracy across optimizers on Model B.}
%     \label{fig:opt_comp}
% \end{figure}

% \subsection{Theoretical Impact of the Momentum Parameter}
% \begin{enumerate}
%     \item \textbf{Damping Oscillations in Anisotropic Ravines}: Loss surfaces of convolutional networks exhibit ill-conditioned curvature (ravines), where curvature is steep along one direction but gentle along another. Standard SGD oscillates perpendicular to the ravine walls because large gradients dominate, impeding forward progress. Momentum incorporates fraction $\gamma$ of prior velocity, canceling out alternating perpendicular oscillations while accelerating descent along the valley floor.
%     \item \textbf{Escaping Saddle Points and Plateaus}: Near zero-gradient saddle points ($\|\nabla L(\theta)\| \approx 0$), standard SGD stalls. The velocity vector $v_t$ maintains kinetic inertia, propelling the parameters across plateaus and escaping shallow local minima.
% \end{enumerate}

% \section{Custom Model Training \& Evaluation [25 Marks]}

% \subsection{Training Protocol}
% Both custom architectures were trained for \textbf{20 epochs} with a batch size of 32 using Adam ($\eta=10^{-3}$), equipped with:
% \begin{itemize}
%     \item \texttt{ModelCheckpoint}: Preserves weights achieving peak validation accuracy.
%     \item \texttt{ReduceLROnPlateau}: Halves the learning rate if validation loss plateaus for 4 consecutive epochs.
%     \item \texttt{CSVLogger}: Records epoch-by-epoch loss and accuracy metrics.
% \end{itemize}

% \begin{figure}[H]
%     \centering
%     \begin{subfigure}[b]{0.48\linewidth}
%         \includegraphics[width=\linewidth]{../results/figures/model_a_training_curves.png}
%         \caption{Model A (Standard CNN)}
%     \end{subfigure}
%     \hfill
%     \begin{subfigure}[b]{0.48\linewidth}
%         \includegraphics[width=\linewidth]{../results/figures/model_b_training_curves.png}
%         \caption{Model B (Separable CNN)}
%     \end{subfigure}
%     \caption{Training and validation loss and accuracy curves (20 epochs).}
%     \label{fig:custom_curves}
% \end{figure}

% \begin{figure}[H]
%     \centering
%     \begin{subfigure}[b]{0.48\linewidth}
%         \includegraphics[width=\linewidth]{../results/figures/model_a_confusion_matrix.png}
%         \caption{Model A Confusion Matrix}
%     \end{subfigure}
%     \hfill
%     \begin{subfigure}[b]{0.48\linewidth}
%         \includegraphics[width=\linewidth]{../results/figures/model_b_confusion_matrix.png}
%         \caption{Model B Confusion Matrix}
%     \end{subfigure}
%     \caption{Normalized test confusion matrices across 15 vegetable classes.}
%     \label{fig:custom_cms}
% \end{figure}

% \subsection{Custom Model Comparison Table}

% \begin{table}[H]
% \centering
% \caption{Performance Comparison between Model A and Model B on Test Set}
% \label{tab:custom_comp}
% \resizebox{\linewidth}{!}{
% \begin{tabular}{lccc}
% \toprule
% \textbf{Metric} & \textbf{Model A (Standard CNN)} & \textbf{Model B (Separable CNN)} & \textbf{Observed Delta / Benefit} \\
% \midrule
% Total Parameters & 260,335 & \textbf{41,391} & \textbf{$-$84.1\% reduction} \\
% Trainable Parameters & 259,631 & \textbf{40,175} & \textbf{$-$84.5\% reduction} ($\le 100\text{k}$) \\
% Disk Size (.keras) & 3,135.53 KB (3.06 MB) & \textbf{628.02 KB (0.61 MB)} & \textbf{$-$80.0\% reduction} \\
% Avg Training Time / Epoch & 54.37 s & \textbf{28.60 s} & \textbf{47.4\% faster per epoch} \\
% Test Accuracy (\%) & 97.57\% & \textbf{97.70\%} & \textbf{+0.13\% (superior generalization)} \\
% Macro Precision (\%) & 97.75\% & 97.75\% & Identical high precision \\
% Macro Recall (\%) & 97.57\% & \textbf{97.70\%} & Balanced class sensitivity \\
% Macro F1-Score (\%) & 97.59\% & \textbf{97.70\%} & Robust balanced performance \\
% Inference Latency (CPU) & 15.32 ms & 25.24 ms & Real-time edge classification \\
% \bottomrule
% \end{tabular}
% }
% \end{table}

% \subsection{Trade-Off Discussion: Standard vs. Depthwise Separable Convolutions}
% \begin{enumerate}
%     \item \textbf{Massive Parameter \& Storage Reduction}: Model B slashes trainable parameters from 259,631 down to 40,175 (an $84.5\%$ reduction), shrinking disk storage from $3.06\text{ MB}$ to $628\text{ KB}$.
%     \item \textbf{Significant Training Acceleration}: Per-epoch training time decreased from $54.37\text{ s}$ to $28.60\text{ s}$ (a $47.4\%$ speedup), empirically verifying the theoretical reduction in MAC operations.
%     \item \textbf{Regularization Benefit}: Model B achieved \textbf{97.70\% test accuracy}, exceeding Model A's 97.57\%. Decoupling spatial and cross-channel filtering prevents overparameterization on low-resolution $64 \times 64$ images.
% \end{enumerate}

% \section{Lightweight State-of-the-Art Networks [20 Marks]}

% We fine-tune two premier mobile architectures using transfer learning from ImageNet under identical data splits:
% \begin{enumerate}
%     \item \textbf{MobileNetV2}: Employs inverted residuals with linear bottlenecks. Channels are expanded by factor $t=6$, spatially filtered via depthwise convolution, and projected back through linear bottlenecks.
%     \item \textbf{EfficientNet-B0}: Utilizes compound scaling, systematically scaling depth, width, and resolution with squeeze-and-excitation (MBConv) blocks.
% \end{enumerate}

% \begin{figure}[H]
%     \centering
%     \begin{subfigure}[b]{0.48\linewidth}
%         \includegraphics[width=\linewidth]{../results/figures/mobilenetv2_training_curves.png}
%         \caption{MobileNetV2 Curves}
%     \end{subfigure}
%     \hfill
%     \begin{subfigure}[b]{0.48\linewidth}
%         \includegraphics[width=\linewidth]{../results/figures/efficientnetb0_training_curves.png}
%         \caption{EfficientNet-B0 Curves}
%     \end{subfigure}
%     \caption{Training and validation curves for fine-tuned SOTA architectures.}
%     \label{fig:sota_curves}
% \end{figure}

% \begin{figure}[H]
%     \centering
%     \begin{subfigure}[b]{0.48\linewidth}
%         \includegraphics[width=\linewidth]{../results/figures/mobilenetv2_confusion_matrix.png}
%         \caption{MobileNetV2 Confusion Matrix}
%     \end{subfigure}
%     \hfill
%     \begin{subfigure}[b]{0.48\linewidth}
%         \includegraphics[width=\linewidth]{../results/figures/efficientnetb0_confusion_matrix.png}
%         \caption{EfficientNet-B0 Confusion Matrix}
%     \end{subfigure}
%     \caption{Test set confusion matrices for fine-tuned SOTA architectures.}
%     \label{fig:sota_cms}
% \end{figure}

% \section{Final Comparison \& Trade-Off Discussion [20 Marks]}

% \subsection{Master Architectural Comparison}

% \begin{table}[H]
% \centering
% \caption{Comprehensive Comparison of Custom and SOTA Architectures}
% \label{tab:master_comp}
% \resizebox{\linewidth}{!}{
% \begin{tabular}{lcccc}
% \toprule
% \textbf{Metric} & \textbf{Model A (Standard)} & \textbf{Model B (Separable)} & \textbf{MobileNetV2 (SOTA)} & \textbf{EfficientNet-B0 (SOTA)} \\
% \midrule
% Design Paradigm & From Scratch & From Scratch ($<100\text{k}$) & Pre-trained Transfer & Pre-trained Transfer \\
% Total Parameters & 260,335 & \textbf{41,391} & 2,277,199 & 4,068,786 \\
% Trainable Parameters & 259,631 & \textbf{40,175} & 2,243,087 & 4,026,763 \\
% Disk Size (.keras) & 3.06 MB & \textbf{0.61 MB (628 KB)} & 26.48 MB & 47.20 MB \\
% Training Time / Epoch & 54.37 s & \textbf{28.60 s} & 166.05 s & 237.16 s \\
% Test Accuracy (\%) & 97.57\% & \textbf{97.70\%} & \textbf{99.60\%} & \textbf{99.53\%} \\
% Macro Precision (\%) & 97.75\% & 97.75\% & 99.60\% & 99.54\% \\
% Macro Recall (\%) & 97.57\% & 97.70\% & 99.60\% & 99.53\% \\
% Macro F1-Score (\%) & 97.59\% & 97.70\% & 99.60\% & 99.53\% \\
% CPU Inference Latency & 15.32 ms & 25.24 ms & 156.17 ms & 278.74 ms \\
% Target Deployment Tier & Single-Board Computer & \textbf{Microcontroller (MCU)} & Edge Gateway / Mobile & Edge Server / Jetson \\
% \bottomrule
% \end{tabular}
% }
% \end{table}

% \subsection{In-Depth Trade-Off Analysis}

% \subsubsection{1. Accuracy vs. Memory Footprint}
% Model B exhibits an exceptional accuracy-to-parameter ratio:
% \begin{equation}
%     \text{Efficiency}_{\text{Model B}} = \frac{97.70\%}{41.4\text{k params}} \approx \mathbf{2.36\% \text{ per 1,000 parameters}}
% \end{equation}
% In comparison, MobileNetV2 yields $0.044\%$ per 1,000 parameters, and EfficientNet-B0 yields $0.024\%$. SOTA models require \textbf{55× and 98× more parameters} to deliver a marginal accuracy gain of $+1.90\%$.

% \subsubsection{2. Hardware Feasibility on Bare Microcontrollers}
% Embedded microcontrollers (e.g., STM32H7, ESP32-S3, Nordic nRF5340) possess strictly limited Flash memory (512 KB to 2 MB) and internal SRAM (256 KB to 512 KB). Deploying MobileNetV2 ($26.5\text{ MB}$) or EfficientNet-B0 ($47.2\text{ MB}$) is physically impossible without external memory interfaces, introducing severe latency and power penalties. Conversely, \textbf{Model B occupies only 628 KB in uncompressed FP32 (and under 45 KB when quantized to INT8)}, executing entirely from on-chip SRAM.

% \subsubsection{3. Computational Cost, Latency, and Battery Life}
% On battery-operated IoT nodes, energy consumption is directly proportional to active cycle time ($E = P \cdot t_{\text{inference}}$). Model B achieves an inference latency of \textbf{25.24 ms} on CPU, allowing the microcontroller to return to deep-sleep low-power states immediately ("race-to-sleep"). Deep models such as MobileNetV2 (156 ms) and EfficientNet-B0 (278 ms) suffer substantial memory access overhead and thermal dissipation on constrained processors.

% \subsection{Deployment Decision Matrix}
% \begin{itemize}
%     \item \textbf{Tier 1 (Microcontrollers / MCUs)}: \textbf{Model B (<100k)} is the definitive choice, delivering 97.70\% accuracy within severe sub-1MB Flash and SRAM envelopes.
%     \item \textbf{Tier 2 (Single-Board Computers / Edge Gateways)}: \textbf{MobileNetV2} is recommended when deploying on Raspberry Pi or NVIDIA Jetson platforms where storage ($>32\text{ MB}$) and power are accessible, maximizing accuracy at 99.60\%.
% \end{itemize}

% \section{Conclusion}
% This study demonstrated that custom architectures leveraging Depthwise Separable Convolutions provide near-SOTA performance ($97.70\%$ vs. $99.60\%$) while reducing parameter counts by $98.2\%$ and memory footprint by $97.7\%$. Hardware-aware design choices—such as $64 \times 64$ sensor downscaling, ReLU6 bounded activations for INT8 quantization, and adaptive Adam optimization—enable deep learning to be practically realized on resource-constrained edge hardware.

% \begin{thebibliography}{99}
% \bibitem{ahmed2021} M. I. Ahmed, S. M. Mamun, and A. U. Z. Asif, ``DCNN-Based Vegetable Image Classification Using Transfer Learning: A Comparative Study,'' \textit{5th International Conference on Computer, Communication and Signal Processing (ICCCSP)}, 2021.
% \bibitem{lecun1998} Y. LeCun, L. Bottou, Y. Bengio, and P. Haffner, ``Gradient-Based Learning Applied to Document Recognition,'' \textit{Proceedings of the IEEE}, vol. 86, no. 11, pp. 2278--2324, 1998.
% \bibitem{chollet2017} F. Chollet, ``Xception: Deep Learning with Depthwise Separable Convolutions,'' \textit{IEEE Conference on Computer Vision and Pattern Recognition (CVPR)}, 2017.
% \bibitem{howard2017} A. G. Howard et al., ``MobileNets: Efficient Convolutional Neural Networks for Mobile Vision Applications,'' \textit{arXiv:1704.04861}, 2017.
% \bibitem{sandler2018} M. Sandler, A. Howard, M. Zhu, A. Zhmoginov, and L.-C. Chen, ``MobileNetV2: Inverted Residuals and Linear Bottlenecks,'' \textit{IEEE Conference on Computer Vision and Pattern Recognition (CVPR)}, 2018.
% \bibitem{tan2019} M. Tan and Q. V. Le, ``EfficientNet: Rethinking Model Scaling for Convolutional Neural Networks,'' \textit{International Conference on Machine Learning (ICML)}, 2019.
% \bibitem{murphy2022} K. P. Murphy, \textit{Probabilistic Machine Learning: An Introduction}, MIT Press, 2022.
% \end{thebibliography}

% \end{document}


\documentclass[11pt,a4paper]{article}

% ---------------- PAGE LAYOUT ----------------
\usepackage[margin=1in]{geometry}
\usepackage[T1]{fontenc}
\usepackage{lmodern}

% ---------------- MATHEMATICS ----------------
\usepackage{amsmath,amssymb,amsfonts}

% ---------------- FIGURES AND TABLES ----------------
\usepackage{graphicx}
\graphicspath{{../results/figures/}{./}}
\usepackage{booktabs}
\usepackage{multirow}
\usepackage{float}
\usepackage{caption}
\usepackage{subcaption}

% ---------------- REFERENCES AND LINKS ----------------
\usepackage{cite}
\usepackage{xcolor}
\usepackage{hyperref}

% ---------------- SECTION FORMATTING ----------------
\usepackage{titlesec}

% ---------------- HEADER AND FOOTER ----------------
\usepackage{fancyhdr}

\hypersetup{
    colorlinks=true,
    linkcolor=blue!70!black,
    citecolor=green!50!black,
    urlcolor=blue!80!black
}

\titleformat{\section}
    {\large\bfseries}{\thesection}{1em}{}

\titleformat{\subsection}
    {\normalsize\bfseries}{\thesubsection}{1em}{}

\pagestyle{fancy}
\fancyhf{}
\fancyhead[L]{Assignment 03}
\fancyfoot[C]{\thepage}
\setlength{\headheight}{14pt}

% ==================================================
\begin{document}

% ---------------- TITLE PAGE ----------------
\begin{titlepage}
    \thispagestyle{empty}
    \centering

    {\LARGE Faculty of Engineering\par}
    \vspace{0.2cm}

    {\large
    Department of Electronic \& Telecommunication Engineering\\
    University of Moratuwa\par}

    \vspace{1cm}

    \includegraphics[width=10cm]{logo.png}

    \vspace{1cm}

    {\LARGE
    \textbf{EN3150 - Pattern Recognition \& Machine Intelligence}\par}

    \vspace{0.3cm}

    {\LARGE \textbf{Assignment 03}\par}

    \vspace{0.5cm}

    {\large
    Resource-Constrained CNN for Edge Image Classification\par}

    \vspace{0.4cm}

    {\normalsize
    Empirical Evaluation of Depthwise Separable Convolutions,
    Optimizer Dynamics, and Lightweight SOTA Transfer Learning\par}

    \vfill

    {\large \textbf{Group Members}\par}
    \vspace{0.3cm}

    {\large
    Thevinu Fernando\\
    Januka Sanujaya\\
    Thisura Samuditha\\
    Chamod Dharmakeerthi\par}

    \vspace{0.5cm}

    {\large
    Department of Electronic \& Telecommunication Engineering\par}
    {\large University of Moratuwa\par}

    \vspace{0.5cm}

    {\large GitHub Repository\par}
    {\small
    \url{https://github.com/thevinufernando/cnn-based-vegetable-classification-model}\par}

    \vspace{0.5cm}

    {\large October 17, 2026\par}

\end{titlepage}

% ---------------- TABLE OF CONTENTS ----------------
\tableofcontents
\thispagestyle{fancy}
\clearpage

% ---------------- ABSTRACT ----------------
\phantomsection
\addcontentsline{toc}{section}{Abstract}

\begin{abstract}
Deploying deep convolutional neural networks (CNNs) on resource-constrained edge hardware (e.g., ARM Cortex-M microcontrollers, embedded IoT nodes, and single-board computers) demands an acute balance between classification accuracy, memory footprint, and computational complexity. In this assignment, we investigate edge-oriented image classification across 15 vegetable categories downscaled to $64 \times 64$ pixels. We design, implement, and benchmark two custom architectures from scratch: \textbf{Model A} (Standard CNN, 259,631 trainable parameters) and \textbf{Model B} (Depthwise Separable CNN, 40,175 trainable parameters, adhering strictly to a sub-100k budget). Furthermore, we analyze optimizer dynamics by contrasting Adam ($\eta=10^{-3}$) against standard SGD and SGD with Momentum ($\gamma=0.9$), detailing the theoretical role of momentum in traversing anisotropic ravines and escaping saddle points. Finally, we fine-tune two mobile state-of-the-art architectures (\textbf{MobileNetV2} and \textbf{EfficientNet-B0}) under identical data splits. Our empirical findings demonstrate that Model B achieves \textbf{97.70\% test accuracy}—outperforming Model A (\textbf{97.57\%}) while reducing parameters by \textbf{84.5\%}, disk footprint by \textbf{80.0\%}, and training time per epoch by \textbf{47.4\%}. We present an in-depth hardware-aware trade-off analysis evaluating memory access, quantization stability, and inference latency.
\end{abstract}

\clearpage

% ==================================================
\section{Resource-Constrained CNN for Edge Image Classification}

\subsection{Dataset Selection \& Low-Resolution Edge Sensor Emulation}

Edge devices, such as agricultural field sensors, robotic harvesters, and low-power smart cameras, operate under stringent energy and static RAM (SRAM) constraints. For this study, we selected the \textbf{Vegetable Image Dataset} (from Kaggle), encompassing 21,000 RGB images across 15 distinct classes: \textit{Bean, Bitter Gourd, Bottle Gourd, Brinjal, Broccoli, Cabbage, Capsicum, Carrot, Cauliflower, Cucumber, Papaya, Potato, Pumpkin, Radish, and Tomato}.

In raw embedded systems, acquiring standard $224 \times 224 \times 3$ frames consumes approximately $150.5\text{ KB}$ per uncompressed image buffer. On microcontrollers featuring only $256\text{ KB}$ to $512\text{ KB}$ of total SRAM, buffering high-resolution images leaves insufficient memory for intermediate feature maps and activation buffers. Consequently, downscaling images to a maximum resolution of \textbf{$64 \times 64$ pixels} compresses the raw uncompressed buffer down to \textbf{$12.288\text{ KB}$} ($64 \times 64 \times 3$ bytes). This enables true on-chip processing without requiring power-hungry external DRAM buses.

\subsection{Data Partitioning \& Leakage Prevention}

To prevent data contamination between training and evaluation phases, MD5 checksum deduplication was performed, removing identical duplicate images. The dataset was partitioned into a strict stratified split:

\begin{itemize}
    \item \textbf{Training Set (70\%)}: 14,996 images
    \item \textbf{Validation Set (15\%)}: 3,000 images
    \item \textbf{Testing Set (15\%)}: 3,000 images
    \item \textbf{Total Unique Images}: 20,996 images
\end{itemize}

This fixed split is deterministically recorded in \texttt{splits.csv} to guarantee absolute reproducibility across all experiments.

% ==================================================
\section{Custom Architecture Design}

\subsection{Theoretical Foundations: Standard vs. Depthwise Separable Convolutions}

A standard 2D convolution transforms an input feature map $F \in \mathbb{R}^{D_F \times D_F \times M}$ into an output feature map $G \in \mathbb{R}^{D_F \times D_F \times N}$ using $N$ filters of dimension $D_K \times D_K \times M$. Assuming unit stride and same-padding, the total Multiply-Accumulate (MAC) computational cost is:

\begin{equation}
    \text{Cost}_{\text{Standard}}
    = D_K^2 M N D_F^2
\end{equation}

In contrast, \textbf{Depthwise Separable Convolutions} (pioneered in MobileNets and Xception) factorize this joint spatial-channel operation into two decoupled stages:

\begin{enumerate}
    \item \textbf{Depthwise Convolution}: Applies a single $D_K \times D_K \times 1$ spatial filter per input channel:

    \begin{equation}
        \text{Cost}_{\text{Depthwise}}
        = D_K^2 M D_F^2
    \end{equation}

    \item \textbf{Pointwise Convolution}: A $1 \times 1$ convolution combining linear projections across all $M$ channels to produce $N$ channels:

    \begin{equation}
        \text{Cost}_{\text{Pointwise}}
        = M N D_F^2
    \end{equation}
\end{enumerate}

Summing both stages, the total computational complexity becomes:

\begin{equation}
    \text{Cost}_{\text{Separable}}
    = D_K^2 M D_F^2 + M N D_F^2
\end{equation}

The relative computational and parameter reduction ratio is given by:

\begin{equation}
    \frac{\text{Cost}_{\text{Separable}}}
    {\text{Cost}_{\text{Standard}}}
    = \frac{D_K^2 M D_F^2 + M N D_F^2}
    {D_K^2 M N D_F^2}
    = \frac{1}{N} + \frac{1}{D_K^2}
\end{equation}

For standard $3 \times 3$ kernels ($D_K=3$) and typical channel depths ($N \geq 32$):

\begin{equation}
    \frac{1}{N}+\frac{1}{9}
    \approx \frac{1}{9}
    \approx 11.1\%
\end{equation}

Depthwise separable convolutions therefore achieve approximately an $8\times$ to $9\times$ reduction in MAC operations, corresponding to roughly $88.9\%$ computational savings when the output channel count is sufficiently large.

\subsection{Architectural Specifications \& Parameter Verification}

\subsubsection{Model A: Standard CNN}

Model A interleaves four standard 2D convolution blocks (with Batch Normalization and ReLU), $2 \times 2$ Max Pooling, Global Average Pooling, and a Dense classification head:

\begin{itemize}
    \item \textbf{Block 1}: Conv2D(32, $3 \times 3$) $\rightarrow$ BN $\rightarrow$ ReLU $\rightarrow$ MaxPool($2 \times 2$) $\rightarrow$ Output: $32 \times 32 \times 32$. Params: 864.
    \item \textbf{Block 2}: Conv2D(64, $3 \times 3$) $\rightarrow$ BN $\rightarrow$ ReLU $\rightarrow$ MaxPool($2 \times 2$) $\rightarrow$ Output: $16 \times 16 \times 64$. Params: 18,432.
    \item \textbf{Block 3}: Conv2D(128, $3 \times 3$) $\rightarrow$ BN $\rightarrow$ ReLU $\rightarrow$ MaxPool($2 \times 2$) $\rightarrow$ Output: $8 \times 8 \times 128$. Params: 73,728.
    \item \textbf{Block 4}: Conv2D(128, $3 \times 3$) $\rightarrow$ BN $\rightarrow$ ReLU $\rightarrow$ MaxPool($2 \times 2$) $\rightarrow$ Output: $4 \times 4 \times 128$. Params: 147,456.
    \item \textbf{Head}: GlobalAveragePooling2D $\rightarrow$ Dense(128, ReLU) $\rightarrow$ Dropout(0.3) $\rightarrow$ Dense(15, Softmax). Params: 18,447.
    \item \textbf{Total Parameters}: \textbf{260,335} (Trainable: \textbf{259,631}).
\end{itemize}

\subsubsection{Model B: Lightweight Separable CNN (Sub-100k Budget)}

Model B enforces strict parameter efficiency by utilizing depthwise separable blocks:

\begin{itemize}
    \item \textbf{Entry Conv}: Conv2D(16, $3 \times 3$) $\rightarrow$ BN $\rightarrow$ ReLU6 $\rightarrow$ MaxPool($2 \times 2$) $\rightarrow$ Output: $32 \times 32 \times 16$. Params: 496.
    \item \textbf{Separable Block 1}: DepthwiseConv2D($3 \times 3$) $\rightarrow$ PointwiseConv2D(32, $1 \times 1$) $\rightarrow$ MaxPool($2 \times 2$). Params: 848.
    \item \textbf{Separable Block 2}: DepthwiseConv2D($3 \times 3$) $\rightarrow$ PointwiseConv2D(64, $1 \times 1$) $\rightarrow$ MaxPool($2 \times 2$). Params: 2,720.
    \item \textbf{Separable Block 3}: DepthwiseConv2D($3 \times 3$) $\rightarrow$ PointwiseConv2D(128, $1 \times 1$) $\rightarrow$ MaxPool($2 \times 2$). Params: 9,536.
    \item \textbf{Separable Block 4}: DepthwiseConv2D($3 \times 3$) $\rightarrow$ PointwiseConv2D(128, $1 \times 1$). Params: 18,560.
    \item \textbf{Head}: GlobalAveragePooling2D $\rightarrow$ Dense(64, ReLU) $\rightarrow$ Dropout(0.2) $\rightarrow$ Dense(15, Softmax). Params: 9,231.
    \item \textbf{Total Parameters}: \textbf{41,391} (Trainable: \textbf{40,175}).
\end{itemize}

\noindent\textbf{Constraint Compliance}: Model B satisfies the assignment criterion of at most 100,000 trainable parameters, with 40,175 trainable parameters, approximately 59.8\% below the limit.

\subsection{Hardware-Aware Activation Function Justification}

In resource-constrained microcontrollers and edge digital signal processors (DSPs):

\begin{enumerate}
    \item \textbf{Elimination of Transcendental Exponentials}: Activations such as Sigmoid ($\sigma(x)=\frac{1}{1+e^{-x}}$) and Tanh require exponential evaluations. These operations can be more expensive on microcontrollers without hardware floating-point or optimized math support.

    \item \textbf{Efficient ReLU Execution}: Standard ReLU, $f(x)=\max(0,x)$, is computationally simple and can often be implemented efficiently using processor instructions or optimized neural-network kernels.

    \item \textbf{Quantization Stability via ReLU6}: In Model B, we employ \textbf{ReLU6}, $f(x)=\min(\max(0,x),6)$. Bounded activations can help control dynamic range during quantization. However, actual INT8 quantization scales and zero-points depend on the quantization scheme and calibration data.
\end{enumerate}

% ==================================================
\section{Optimizer Selection \& Tuning}

\subsection{Selected Optimizer: Adam}

We selected \textbf{Adam (Adaptive Moment Estimation)} with an initial learning rate $\eta=0.001$, $\beta_1=0.9$, $\beta_2=0.999$, and $\epsilon=10^{-7}$.

\noindent\textbf{Justification}: Adam computes exponentially decaying moving averages of past gradients ($m_t$, first moment) and squared gradients ($v_t$, second moment). These estimates adapt parameter updates to the scale of recent gradients and can be useful when training deep vision architectures.

\subsection{Comparative Analysis: Standard SGD vs. SGD with Momentum}

Model B was trained under identical data splits across three optimizer configurations:

\begin{enumerate}
    \item \textbf{Adam}: $\eta=0.001$.
    \item \textbf{Standard SGD}: $\theta_{t+1}=\theta_t-\eta\nabla L(\theta_t)$, with $\eta=0.01$ and momentum $=0.0$.
    \item \textbf{SGD with Momentum}: $v_{t+1}=\gamma v_t+\eta\nabla L(\theta_t)$, $\theta_{t+1}=\theta_t-v_{t+1}$, with $\gamma=0.9$.
\end{enumerate}

\begin{figure}[H]
    \centering
    \includegraphics[width=0.92\linewidth]{optimizer_comparison.png}
    \caption{Loss convergence and validation accuracy across optimizers on Model B.}
    \label{fig:opt_comp}
\end{figure}

\subsection{Theoretical Impact of the Momentum Parameter}

\begin{enumerate}
    \item \textbf{Damping Oscillations in Anisotropic Ravines}: Loss surfaces of convolutional networks can exhibit ill-conditioned curvature. Momentum accumulates a running velocity that may reduce oscillations and accelerate progress along consistent gradient directions.

    \item \textbf{Behavior Near Saddle Points and Plateaus}: Momentum can carry parameters through regions with small gradients if a nonzero velocity has accumulated. However, escaping a saddle point is not guaranteed and depends on the loss geometry and gradient noise.
\end{enumerate}

% ==================================================
\section{Custom Model Training \& Evaluation}

\subsection{Training Protocol}

Both custom architectures were trained for \textbf{20 epochs} with a batch size of 32 using Adam ($\eta=10^{-3}$), equipped with:

\begin{itemize}
    \item \texttt{ModelCheckpoint}: Preserves weights achieving peak validation accuracy.
    \item \texttt{ReduceLROnPlateau}: Halves the learning rate if validation loss plateaus for 4 consecutive epochs.
    \item \texttt{CSVLogger}: Records epoch-by-epoch loss and accuracy metrics.
\end{itemize}

\begin{figure}[H]
    \centering
    \begin{subfigure}[b]{0.48\linewidth}
        \includegraphics[width=\linewidth]{model_a_training_curves.png}
        \caption{Model A (Standard CNN)}
    \end{subfigure}
    \hfill
    \begin{subfigure}[b]{0.48\linewidth}
        \includegraphics[width=\linewidth]{model_b_training_curves.png}
        \caption{Model B (Separable CNN)}
    \end{subfigure}
    \caption{Training and validation loss and accuracy curves (20 epochs).}
    \label{fig:custom_curves}
\end{figure}

\begin{figure}[H]
    \centering
    \begin{subfigure}[b]{0.48\linewidth}
        \includegraphics[width=\linewidth]{model_a_confusion_matrix.png}
        \caption{Model A Confusion Matrix}
    \end{subfigure}
    \hfill
    \begin{subfigure}[b]{0.48\linewidth}
        \includegraphics[width=\linewidth]{model_b_confusion_matrix.png}
        \caption{Model B Confusion Matrix}
    \end{subfigure}
    \caption{Normalized test confusion matrices across 15 vegetable classes.}
    \label{fig:custom_cms}
\end{figure}

\subsection{Custom Model Comparison Table}

\begin{table}[H]
\centering
\caption{Performance Comparison between Model A and Model B on Test Set}
\label{tab:custom_comp}
\resizebox{\linewidth}{!}{
\begin{tabular}{lccc}
\toprule
\textbf{Metric} & \textbf{Model A (Standard CNN)} & \textbf{Model B (Separable CNN)} & \textbf{Observed Delta / Benefit} \\
\midrule
Total Parameters & 260,335 & \textbf{41,391} & \textbf{$-$84.1\% reduction} \\
Trainable Parameters & 259,631 & \textbf{40,175} & \textbf{$-$84.5\% reduction} ($\leq 100\text{k}$) \\
Disk Size (.keras) & 3,135.53 KB (3.06 MB) & \textbf{628.02 KB (0.61 MB)} & \textbf{$-$80.0\% reduction} \\
Avg Training Time / Epoch & 54.37 s & \textbf{28.60 s} & \textbf{47.4\% faster per epoch} \\
Test Accuracy (\%) & 97.57\% & \textbf{97.70\%} & \textbf{+0.13 percentage points} \\
Macro Precision (\%) & 97.75\% & 97.75\% & Identical reported precision \\
Macro Recall (\%) & 97.57\% & \textbf{97.70\%} & Balanced class sensitivity \\
Macro F1-Score (\%) & 97.59\% & \textbf{97.70\%} & Balanced performance \\
Inference Latency (CPU) & 15.32 ms & 25.24 ms & See latency discussion \\
\bottomrule
\end{tabular}
}
\end{table}

\subsection{Trade-Off Discussion: Standard vs. Depthwise Separable Convolutions}

\begin{enumerate}
    \item \textbf{Parameter and Storage Reduction}: Model B reduces trainable parameters from 259,631 to 40,175, an 84.5\% reduction, and reduces the reported model file size from 3.06 MB to approximately 628 KB.

    \item \textbf{Training Acceleration}: Per-epoch training time decreased from 54.37 s to 28.60 s, a reported speedup of approximately 47.4\%.

    \item \textbf{Regularization Benefit}: Model B achieved 97.70\% test accuracy compared with Model A's 97.57\%. This is an observed difference; further repeated runs would be needed to establish whether it is statistically reliable.
\end{enumerate}

% ==================================================
\section{Lightweight State-of-the-Art Networks}

We fine-tune two mobile architectures using transfer learning from ImageNet under identical data splits:

\begin{enumerate}
    \item \textbf{MobileNetV2}: Employs inverted residuals with linear bottlenecks. Channels are expanded by a factor $t=6$, spatially filtered via depthwise convolution, and projected back through linear bottlenecks.

    \item \textbf{EfficientNet-B0}: Uses compound scaling and MBConv blocks with squeeze-and-excitation modules.
\end{enumerate}

\begin{figure}[H]
    \centering
    \begin{subfigure}[b]{0.48\linewidth}
        \includegraphics[width=\linewidth]{mobilenetv2_training_curves.png}
        \caption{MobileNetV2 Curves}
    \end{subfigure}
    \hfill
    \begin{subfigure}[b]{0.48\linewidth}
        \includegraphics[width=\linewidth]{efficientnetb0_training_curves.png}
        \caption{EfficientNet-B0 Curves}
    \end{subfigure}
    \caption{Training and validation curves for fine-tuned SOTA architectures.}
    \label{fig:sota_curves}
\end{figure}

\begin{figure}[H]
    \centering
    \begin{subfigure}[b]{0.48\linewidth}
        \includegraphics[width=\linewidth]{mobilenetv2_confusion_matrix.png}
        \caption{MobileNetV2 Confusion Matrix}
    \end{subfigure}
    \hfill
    \begin{subfigure}[b]{0.48\linewidth}
        \includegraphics[width=\linewidth]{efficientnetb0_confusion_matrix.png}
        \caption{EfficientNet-B0 Confusion Matrix}
    \end{subfigure}
    \caption{Test set confusion matrices for fine-tuned SOTA architectures.}
    \label{fig:sota_cms}
\end{figure}

% ==================================================
\section{Final Comparison \& Trade-Off Discussion}

\subsection{Master Architectural Comparison}

\begin{table}[H]
\centering
\caption{Comprehensive Comparison of Custom and SOTA Architectures}
\label{tab:master_comp}
\resizebox{\linewidth}{!}{
\begin{tabular}{lcccc}
\toprule
\textbf{Metric} & \textbf{Model A (Standard)} & \textbf{Model B (Separable)} & \textbf{MobileNetV2 (SOTA)} & \textbf{EfficientNet-B0 (SOTA)} \\
\midrule
Design Paradigm & From Scratch & From Scratch ($<100\text{k}$) & Pre-trained Transfer & Pre-trained Transfer \\
Total Parameters & 260,335 & \textbf{41,391} & 2,277,199 & 4,068,786 \\
Trainable Parameters & 259,631 & \textbf{40,175} & 2,243,087 & 4,026,763 \\
Disk Size (.keras) & 3.06 MB & \textbf{0.61 MB} & 26.48 MB & 47.20 MB \\
Training Time / Epoch & 54.37 s & \textbf{28.60 s} & 166.05 s & 237.16 s \\
Test Accuracy (\%) & 97.57\% & \textbf{97.70\%} & \textbf{99.60\%} & \textbf{99.53\%} \\
Macro Precision (\%) & 97.75\% & 97.75\% & 99.60\% & 99.54\% \\
Macro Recall (\%) & 97.57\% & 97.70\% & 99.60\% & 99.53\% \\
Macro F1-Score (\%) & 97.59\% & 97.70\% & 99.60\% & 99.53\% \\
CPU Inference Latency & 15.32 ms & 25.24 ms & 156.17 ms & 278.74 ms \\
Target Deployment Tier & Single-Board Computer & \textbf{Small Edge Device} & Edge Gateway / Mobile & Edge Server / Jetson \\
\bottomrule
\end{tabular}
}
\end{table}

\subsection{In-Depth Trade-Off Analysis}

\subsubsection{1. Accuracy vs. Memory Footprint}

Model B exhibits a high accuracy-to-parameter ratio:

\begin{equation}
    \text{Efficiency}_{\text{Model B}}
    = \frac{97.70}{41.4}
    \approx 2.36
\end{equation}

This corresponds to approximately 2.36 percentage points of test accuracy per 1,000 trainable parameters. Such a ratio is a descriptive metric rather than a standardized measure of model quality.

\subsubsection{2. Hardware Feasibility on Bare Microcontrollers}

Embedded microcontrollers have limited Flash memory and SRAM. The reported MobileNetV2 and EfficientNet-B0 model file sizes exceed the on-chip storage of many small microcontrollers. Actual feasibility depends on model conversion, quantization, Flash capacity, activation memory, operator support, and inference runtime.

Model B has 40,175 trainable parameters. At one byte per parameter, its weights alone would require approximately 40 KB when successfully quantized to INT8. This does not include all model metadata, non-trainable parameters, intermediate activation buffers, or runtime memory.

\subsubsection{3. Computational Cost, Latency, and Battery Life}

On battery-operated IoT nodes, inference energy can be approximated as $E=P\cdot t$ when average power $P$ is known. Model B has a reported CPU inference latency of 25.24 ms. Actual energy consumption and suitability for a particular microcontroller should be measured on the target hardware.

\subsection{Deployment Decision Matrix}

\begin{itemize}
    \item \textbf{Tier 1 (Resource-Constrained Edge Devices)}: Model B is the strongest candidate when the trainable parameter budget is below 100,000, subject to validation of memory use and operator support on the target hardware.

    \item \textbf{Tier 2 (Single-Board Computers / Edge Gateways)}: MobileNetV2 is a candidate when additional memory and compute are available and its measured accuracy and latency meet the deployment requirements.
\end{itemize}

% ==================================================
\section{Conclusion}

This study compared custom CNN architectures using standard and depthwise separable convolutions with two pretrained mobile architectures. The reported results show that Model B achieves 97.70\% test accuracy with 40,175 trainable parameters, while MobileNetV2 and EfficientNet-B0 achieve higher reported accuracy at a substantially greater model size.

Hardware-aware choices, including $64 \times 64$ input images, bounded activations, and lightweight convolutional blocks, can improve the practicality of image classification on resource-constrained edge devices. Final deployment suitability should be established using quantized model size, peak runtime memory, and latency measurements on the intended hardware.

% ==================================================
% REFERENCES
% ==================================================
\clearpage
\phantomsection
\addcontentsline{toc}{section}{References}

\begin{thebibliography}{99}

\bibitem{ahmed2021}
M. I. Ahmed, S. M. Mamun, and A. U. Z. Asif,
``DCNN-Based Vegetable Image Classification Using Transfer Learning: A Comparative Study,''
\textit{5th International Conference on Computer, Communication and Signal Processing (ICCCSP)}, 2021.

\bibitem{lecun1998}
Y. LeCun, L. Bottou, Y. Bengio, and P. Haffner,
``Gradient-Based Learning Applied to Document Recognition,''
\textit{Proceedings of the IEEE}, vol. 86, no. 11, pp. 2278--2324, 1998.

\bibitem{chollet2017}
F. Chollet,
``Xception: Deep Learning with Depthwise Separable Convolutions,''
\textit{IEEE Conference on Computer Vision and Pattern Recognition (CVPR)}, 2017.

\bibitem{howard2017}
A. G. Howard et al.,
``MobileNets: Efficient Convolutional Neural Networks for Mobile Vision Applications,''
\textit{arXiv:1704.04861}, 2017.

\bibitem{sandler2018}
M. Sandler, A. Howard, M. Zhu, A. Zhmoginov, and L.-C. Chen,
``MobileNetV2: Inverted Residuals and Linear Bottlenecks,''
\textit{IEEE Conference on Computer Vision and Pattern Recognition (CVPR)}, 2018.

\bibitem{tan2019}
M. Tan and Q. V. Le,
``EfficientNet: Rethinking Model Scaling for Convolutional Neural Networks,''
\textit{International Conference on Machine Learning (ICML)}, 2019.

\bibitem{murphy2022}
K. P. Murphy,
\textit{Probabilistic Machine Learning: An Introduction},
MIT Press, 2022.

\end{thebibliography}

\end{document}
